\section{Conclusão}

Nesse trabalho, conseguimos determinar com alta precisão o valor do comprimento de onda do laser de HeNe e o índice de refração do ar em CNTP. Também conseguimos encontrar um valor para o índice de refração do vidro que é condizente com os valores conhecidos para outros vidros. Concluímos então que a interferometria é um método eficaz para fazer medidas precisas. Em particular, o interferômetro de Michelson é uma montagem relativamente simples e pode ser usada para fazer diversas medidas ópticas. 
